\capitulosinnumero{Justificación y objetivos}

Para comprender la necesidad y los objetivos de este Trabajo de Fin de Grado, es necesario presentar mi situación personal y la motivación que me ha llevado a querer desarrollar un videojuego como este. En primer lugar, debo mencionar que soy una persona transgénero y llevo conviviendo con esa etiqueta durante casi la mitad de mi vida. Algo tan predominante en mi día a día, como lo es mi identidad de género, ha sido un tema que he tenido que aprender a gestionar y comprender a lo largo de los años. La sociedad en la que vivimos no siempre es comprensiva ni respetuosa con las personas trans, y esto ha generado en mí una serie de experiencias y emociones que han influido tanto en mi desarrollo personal como en el académico.

En segundo lugar, desde pequeño siempre quise crear videojuegos gracias a la influencia de mis padres, quienes me introdujeron en este medio a través de obras de todo tipo. En cuanto a la parte de programación, siempre me había llamado la atención aprender a programar, pero no tuve la oportunidad de centrarme en ello hasta que entré en la universidad. Al llegar al instituto, no tenía nada más en mente que entender mi identidad de género y sexual, ya que no me sentía cómodo con mi cuerpo y no comprendía el porqué. En la educación actual en España apenas se informa sobre estas realidades y mucho menos sobre cómo gestionarlas. Si hubiera contado con el apoyo o la información adecuada en aquel momento, probablemente habría podido dar ese paso hacia la programación y otros intereses mucho antes, además de haber disfrutado más de mi infancia y adolescencia.

Por ello, con este proyecto busco unir mi pasión por los videojuegos y la programación con mi vivencia personal como persona transgénero, con la intención de crear un título que no solo entretenga, sino que también trate de cerca la identidad de género y la salud mental. A través de este trabajo espero aportar un grano de arena a la visibilización de estas cuestiones y ofrecer una experiencia interactiva con la que puedan sentirse identificados otros jugadores que estén pasando por situaciones parecidas. Cabe destacar que este proyecto no pretende hablar en nombre de toda la comunidad transgénero, sino representar de manera respetuosa y sincera una parte de estas vivencias sin intentar abarcar todas las realidades posibles. Aunque existen algunos videojuegos que tocan esta temática, siguen siendo pocos y rara vez se centran en el proceso interno de descubrir la propia identidad y el desgaste que conlleva aceptarla frente al entorno.

Partiendo de esta motivación, el objetivo general del proyecto consiste en desarrollar un prototipo funcional de videojuego narrativo en \ac{2D} jugado en primera persona (con sensación de profundidad 2.5D) que explore una crisis de identidad desde la perspectiva de género mediante mecánicas adaptativas. Este objetivo principal se divide en la redacción del Documento de Diseño de Juego (\ac{GDD}) y su dirección artística contando con asesoramiento psicológico, la programación en segundo plano de un Director de Alucinación que evalúe la forma de jugar del usuario, la creación de una base de código flexible con \textit{ScriptableObjects} para aplicar «Mecanismos de Defensa» y cambios en la interfaz (\ac{UI}) pidiendo siempre el consentimiento del jugador, y el desarrollo de una demo jugable en Unity.